\section{Implementation}
\label{sec:impl}

\polyforge{} is a working system, not a simulator with a paper around it.
This section records what was built, the gates it is held to, and the two
defects that wiring the live evaluation arm surfaced before any cluster
time was spent.

\paragraph{Components} The control plane (Go; \evfile{cmd/control-plane},
\evfile{internal/}) exposes tenant, project and API-key APIs over
PostgreSQL with row-level security (an SQLite path exists for
development), rate limiting, telemetry ingestion, and canary and
isolation-ladder admin APIs. The AI gateway (Go; \evfile{cmd/ai-gateway},
\evfile{internal/ai}) routes chat across backends by plan, length and
budget, and hosts the semantic cache, vector search and the agent loop.
The classifier (\evfile{internal/classifier}, trained by
\evfile{cmd/classifier-train}) turns 12-feature windows into taxonomy
labels every 10\,s with PSI drift detection. The \jcac{} planner
(\evfile{services/planner}) imports the simulator's controller from
\evfile{research/jcac\_sim}, so offline and online planning are one code
path. The operator (Go; \evfile{cmd/operator}, Helm chart under
\evfile{deploy/helm}) reconciles the four custom resources with the
guardrails of Section~\ref{sec:design-loop}, tears down in finalizer
order, and traces every reconcile. The evaluation harness (\evfile{eval/})
is YAML-driven, carries the tuned baselines, writes DuckDB results, and
has a kind-cluster live backend.

\paragraph{Engineering gates} Every change passes the same battery before
commit: \texttt{gofmt}, \texttt{go vet}, \texttt{go build}, \texttt{go
test ./... -count=1} across 26 Go packages, the Python suites for the
harness, planner and analysis, and OpenAPI linting. Continuous integration
additionally runs the race detector, a coverage gate, golangci-lint and the
row-level-security suite against a live PostgreSQL service. Tests that only
CI can execute are reported as unproven until a green CI run, and the
session log records how every change was verified. Security-relevant
details received the treatment of research claims: API-key identifiers are
drawn independently of the secret after an audit found the identifier
leaked 64 of 192 entropy bits; keys are hashed with SHA-256 rather than
argon2id, sound because the secrets are 192-bit random strings, and
recorded as a deliberate deviation precisely because it invites the
question; a recovery middleware returns structured errors without stack
details; the authentication packages are held at roughly 95\% coverage.

\paragraph{Deviations from the original plan} Three tools named in the
original roadmap were replaced with equivalent outcomes and the
replacement recorded rather than silently made: the chi router by Go
1.22's standard-library method patterns; sqlc and golang-migrate by
hand-written SQL behind repository interfaces plus an idempotent migrator
(type-safe access and a versioned schema still hold, and row-level
security is proven by test); argon2id by SHA-256 as above. The evaluation
reports latency at p95 rather than the roadmap's p99, for a reason
Section~\ref{sec:method-metrics} gives.

\subsection{The live arm, and two defects found at the desk}
\label{sec:impl-live}

The live campaigns (Section~\ref{sec:eval-live}) run on a kind cluster
with harness-built images, a k6 load generator, real HPA for the baseline
arm, and the operator plus planner for the \jcac{} arm. Wiring that arm
surfaced two defects, each of which would have silently invalidated a
cluster session.

The first was demand plumbing that could never have authenticated. The
operator's feature-demand source sent its token as a bearer, which the
control plane treats strictly as a JWT; every read would have returned
401, and because the \jcac{} loop degrades by design into silent fallback,
the arm would have run ``live'' while planning nothing. The features
endpoint---that endpoint only---now accepts the platform admin key, and
the change is unit-tested in both directions: the admin key reads any
tenant's features and still cannot touch other tenant routes.

The second was last-writer-wins actuation on a shared target. All eight
evaluation tenants' Policies name one shared Deployment, and each
reconcile overwrote the others' replica counts, which would have pinned
the target at a single tenant's share---the joint controller crippled by
wiring and labelled as the joint controller. Policies that resolve to the
same target now sum their clamped contributions, deletion re-enqueues the
siblings, and the behaviour is unit-tested.

Both fixes travel with a more general honesty mechanism. The harness ends
operator installation with \texttt{kubectl wait --for=condition=Applied}
on every Policy, so a run in which the operator never actuated fails
before any load is generated; ``the controller was actually in the loop''
is an executable gate rather than a promise. Capacity parity is a cell
property under test: every arm starts at the same replica count and is
ceilinged by the same cluster limits, so no arm can win by being handed
more cluster.

\subsection{Scale behaviour of the deployed planner}
\label{sec:impl-scaling}

The deployed \texttt{PlannerCore.plan} path was driven from 8 to 256
tenants on one laptop core. The fitted growth exponent is 1.45: the joint
fairness term couples tenants, and super-linearity is the measured price of
jointness. The p95 cycle time first exceeds the operator's 3\,s request
timeout at 128 tenants and the 10\,s control period at 256. Past the
timeout the failure mode is the designed one---the operator holds the last
good plan and reports fallback status. The levers past the crossover were
then measured rather than named: hash-based planning cells hold the
per-cell p95 under the timeout through 1,024 tenants, at a global-Jain
cost the first cell assignment exceeded and a de-aliased assignment brings
within bound; the wall reproduces at the same width through the shipped Go
operator (Section~\ref{sec:threats}). Coordination costs outside the
planner---custom-resource write fan-out, telemetry aggregation---are not
covered by this microbenchmark.
